\documentclass[11pt,a4paper]{article}
\usepackage[margin=1in]{geometry}
\usepackage{amsmath,amssymb,bm}
\usepackage{mathtools}
\usepackage{enumitem}
\usepackage{hyperref}
\setlist[itemize]{leftmargin=1.25em}
\setlength{\parskip}{0.4em}
\setlength{\parindent}{0pt}

\title{Sampling-Stage Dynamics Guidance in Collision\_sim\\Mathematical Summary}
\author{Code-level derivation from \texttt{models/diffusion.py} and \texttt{models/autoencoder.py}}
\date{\today}

\begin{document}
\maketitle

\section{Scope}
This document summarizes how dynamics guidance is applied \emph{during sampling} in the current codebase.
It does not describe training-time parameter updates.

\section{Base Diffusion Sampling Update}
At reverse step $t$, the base update is:
\begin{align}
\text{DDPM:}\quad
\bm{x}_{t-1}
&=
c_0 \left(\bm{x}_t - c_1\,\epsilon_\theta(\bm{x}_t,t,\bm{c})\right)
 + \sigma_t \bm{z},
\\
\text{DDIM:}\quad
\bm{x}_{t-1}
&=
\sqrt{\bar{\alpha}_{t-1}}\,\hat{\bm{x}}_0
 + \sqrt{1-\bar{\alpha}_{t-1}}\,\epsilon_\theta(\bm{x}_t,t,\bm{c}),
\\
\hat{\bm{x}}_0
&=
\frac{\bm{x}_t-\sqrt{1-\bar{\alpha}_t}\,\epsilon_\theta(\bm{x}_t,t,\bm{c})}
{\sqrt{\bar{\alpha}_t}}.
\end{align}

If guidance is enabled and $t \le t_{\text{start}}$, an additional guidance step is applied:
\begin{align}
t_{\text{start}} &= \max\{1,\operatorname{round}(T(1-r_{\text{start}}))\},
\\
\bm{x}_{t-1} &\leftarrow \operatorname{ApplyGuidance}(\bm{x}_{t-1}).
\end{align}

\section{Guidance Update in Sample Space}
\subsection{Standardized-to-physical velocity}
The sampled state is velocity in standardized space. The code maps it to physical velocity:
\begin{align}
\bm{v} = \bm{x}\odot \bm{\sigma}_v,
\end{align}
where $\bm{\sigma}_v$ is the velocity standard deviation.

\subsection{Gradient step}
Let the total objective be
\begin{align}
J(\bm{v}) = J_{\text{dyn}}(\bm{v}) + J_{\text{col}}(\bm{v}).
\end{align}
Then each inner guidance iteration performs
\begin{align}
\bm{x}\leftarrow \bm{x}-\lambda \nabla_{\bm{x}}J(\bm{v}),
\end{align}
with optional gradient-norm clipping.

\section{Dynamics Objective $J_{\text{dyn}}$}
Define
\begin{align}
\operatorname{MM}(q,m) &= \frac{\sum q\,m}{\sum m+\varepsilon},
&
\phi(q;q_{\max},m) &= \operatorname{MM}(\operatorname{ReLU}(q-q_{\max})^2,m).
\end{align}

From sampled velocity sequence $\{\bm{v}_t\}_{t=1}^{H}$:
\begin{align}
\bm{a}_t &= \frac{\bm{v}_t-\bm{v}_{t-1}}{\Delta t},
&
\bm{j}_t &= \frac{\bm{a}_t-\bm{a}_{t-1}}{\Delta t},
\\
\theta_t &= \operatorname{atan2}(v_{t,y},v_{t,x}),
&
r_t &= \frac{|\operatorname{wrap}(\theta_t-\theta_{t-1})|}{\Delta t},
\\
\bar{s}_t &= \tfrac{1}{2}(\|\bm{v}_t\|+\|\bm{v}_{t-1}\|),
&
\kappa_t &= \frac{r_t}{\max(\bar{s}_t,s_{\min})},
\\
a^{\text{lat}}_t &= r_t\,\bar{s}_t.
\end{align}

The objective is a weighted sum of soft-limit penalties:
\begin{align}
J_{\text{dyn}}
&=
w_a\,\phi(\|\bm{a}\|;a_{\max},m_{\text{move}})
+w_j\,\phi(\|\bm{j}\|;j_{\max},m_{\text{jerk}})
+w_r\,\operatorname{MM}(\operatorname{ReLU}(r-r_{\lim})^2,\mathbf{1})
\nonumber\\
&\quad
+w_\kappa\,\phi(\kappa;\kappa_{\max},m_{\text{move}})
+w_{\text{lat}}\,\phi(a^{\text{lat}};a^{\text{lat}}_{\max},m_{\text{move}})
+w_{\text{slip}}\,\phi(\rho;\rho_{\max},m_{\text{move}})
\nonumber\\
&\quad
+w_{\text{rev}}\,\operatorname{MM}\!\left(\operatorname{ReLU}(-( \bm{v}\!\cdot\!\bm{d}_{\text{travel}}+\tau_{\text{rev}}))^2,m_{\text{speed}}\right)
+w_{\text{low}}\,\operatorname{MM}(r^2,m_{\text{low}}).
\end{align}

If bicycle-curvature mode is enabled:
\begin{align}
\kappa_{\text{bicycle}} = \frac{\tan(\delta_{\max})}{L},\qquad
\kappa_{\lim} = \min(\kappa_{\text{cfg}},\kappa_{\text{bicycle}}).
\end{align}

\section{Collision-Mode Objective $J_{\text{col}}$ (Optional)}
Adversary rollout from sampled velocity:
\begin{align}
\bm{p}^{\text{adv}}_t
=
\bm{p}^{\text{adv}}_0+\sum_{\tau=1}^{t}\bm{v}^{\text{adv}}_{\tau}\Delta t.
\end{align}

Target contact point from mode index (11--53):
\begin{align}
\bm{p}^{\text{contact}}_t
=
\bm{p}^{\text{target}}_t
+R(\psi^{\text{target}}_t)\Delta\bm{p}_{\text{mode}}(L,W),\quad
d_t^2=\|\bm{p}^{\text{adv}}_t-\bm{p}^{\text{contact}}_t\|^2.
\end{align}

With softmin over time:
\begin{align}
\alpha_t = \operatorname{softmax}\!\left(-\frac{d_t^2}{\tau}\right),\quad
\tilde{d}^2=\sum_t \alpha_t d_t^2,\quad
J_{\text{col}} = w_{\text{col}}\mathbb{E}\!\left[\frac{\tilde{d}^2}{\sigma^2}\right].
\end{align}
Optional approach and closing-speed terms are added in code.

\section{Algorithm Summary}
\begin{enumerate}
\item Compute base reverse update (DDPM or DDIM).
\item If guidance condition is met, convert to physical velocity.
\item Evaluate $J_{\text{dyn}}(+J_{\text{col}})$.
\item Update sample by gradient descent in sample space.
\item Continue reverse loop to $t=0$.
\end{enumerate}

\section{Interpretation}
Dynamics guidance is \textbf{inference-time control} on sampled trajectories.
It biases generated trajectories without directly changing model weights in that reverse step.

\end{document}
